\documentclass{situationist-preprint}

\title{McTaggart's Argument Inverted}
\msid{TSM-01-b}
\msfamily{TSM-01}
\msdiscipline{Time, Space and Modality}
\mstype{Companion manuscript}
\mskeywords{McTaggart; A-series; B-series; change; unreality of time; modus tollens}

\begin{document}
\maketitle

\begin{abstract}
McTaggart argued that time requires change, that change requires the A-series, and that the A-series is contradictory, so that time is unreal. TSM-01 accepts the third premise and rejects the conclusion: the A-series is incoherent, and time is real and constituted by the B-series. We show that this is not a selective acceptance of McTaggart's premises but the argument run in the other direction. Its first two premises are plausible only in different senses of \enquote{change}, and the single premise the argument needs, that time requires change of A-properties, is the one that the reality of time refutes. Inverted, McTaggart's argument is an argument for the B-theory.
\end{abstract}

\section{The argument and its inversion}

\citet{McTaggart1908} distinguishes two series in which positions in time are ordered: the A-series of past, present and future, and the B-series of earlier and later. His argument has three premises.

\begin{assumption}[Time requires change]\label{a:m1}
There is no time without change.
\end{assumption}

\begin{assumption}[Change requires the A-series]\label{a:m2}
The only change there could be is change in A-determinations: an event's ceasing to be future, becoming present and becoming past.
\end{assumption}

\begin{assumption}[Incoherence]\label{a:m3}
The A-series is incoherent: every event would have to be past, present and future, which are incompatible, and the attempt to remove the contradiction by tensing the copula generates a regress in which it recurs.
\end{assumption}

From the three, time is unreal. TSM-01 accepts Assumption~\ref{a:m3}, holds that time is real, and holds that time is constituted by the B-series, which is the causal and metrical structure of spacetime. Since the argument is valid, the combination requires that Assumption~\ref{a:m1} or Assumption~\ref{a:m2} be false. The question is which, and whether the choice is principled.

\begin{proposition}\label{p:equiv}
The argument needs only the conjunction of Assumptions~\ref{a:m1} and~\ref{a:m2}, which is equivalent to the claim that there is no time without a changing A-series.
\end{proposition}

\begin{proof}
If there is no time without change, and every change is a change in A-determinations, then there is no time without a changing A-series; conversely, the latter claim, with Assumption~\ref{a:m3}, yields the conclusion without either premise separately.
\end{proof}

Call the conjunction the \emph{bridge}. The argument has the form: bridge, incoherence, so no time. Its inversion has the form: time, incoherence, so not the bridge. Which direction is rational depends on comparative credibility. Assumption~\ref{a:m3} is common to both. The question is whether the bridge is more credible than the reality of time.

\section{The bridge equivocates}

Each conjunct of the bridge is plausible only in its own sense of \enquote{change}. Distinguish \emph{B-change}, an object's having one property at one time and an incompatible property at a later time \citep[§442]{Russell1903}, from \emph{A-change}, an event's passing from future to present to past.

\begin{proposition}\label{p:split}
Assumption~\ref{a:m1} is credible only for B-change, and Assumption~\ref{a:m2} only for A-change.
\end{proposition}

\begin{proof}
For Assumption~\ref{a:m1}: the considerations that recommend it are that a time at which nothing differs from any other time seems idle, and that temporal order is known through the succession of states. Both concern variation of things across times, which is B-change; neither mentions events' changing their A-determinations. For Assumption~\ref{a:m2}: McTaggart's ground is that a poker hot on Monday and cold on Tuesday does not change, since it is always true that it is hot on Monday; the only facts about an event that differ at different times are its A-determinations. This ground presupposes that change must be a change in the truth-value of a fixed proposition over time, which is A-change by definition. Read as a claim about B-change, Assumption~\ref{a:m2} is false: the poker is hot at one time and cold at a later one, which is all that B-change is. Read as a claim about A-change, Assumption~\ref{a:m1} has no independent support.
\end{proof}

By Propositions~\ref{p:equiv} and~\ref{p:split}, the bridge has no reading on which both conjuncts are credible. Its credibility derives from the equivocation, and the reality of time does not. Our evidence for temporal order and duration is the evidence for the causal and metrical structure of spacetime, which is as good as any evidence in physics. It is therefore the bridge, not the reality of time, that the incoherence of the A-series refutes. What survives is that time is not the A-series and needs no A-change, which is the claim that time is the B-series. Assumption~\ref{a:m1} in its B-reading is not needed for this; whether it is true is the separate question of TSM-17, which argues that time without change is possible.

\section{Objections and replies}

\begin{objection}[The C-series]
McTaggart replies that without the A-series the B-series reduces to an order without direction, a C-series, and an undirected order is not time.
\end{objection}

\begin{reply}
The B-relation is asymmetric, and the asymmetry is grounded in physical structure: in the causal structure of spacetime, which separates the two halves of each light cone, and in the thermodynamic asymmetry, which tells them apart and from which the asymmetries of records and influence derive (TSM-05). An order with a physically grounded direction is not a C-series. McTaggart's assumption that only the A-series could direct the order is a further instance of the bridge.
\end{reply}

\begin{objection}[Dummett]
\citet{Dummett1960} defends McTaggart on the ground that a complete description of reality cannot be given from no temporal standpoint, since the use of tensed and indexical expressions is ineliminable; reality is therefore either incoherent or not complete in any description, and either way time is not as the B-theory says.
\end{objection}

\begin{reply}
The ineliminability is that of indexical thought, which holds equally of \enquote{here} and \enquote{I}. A complete description of the world need not contain the self-locating information a thinker has about her place in it; EPI-18-c and REL-19-b treat \enquote{now} in just this way, as carrying self-locating rather than tensed information. What Dummett's point establishes is that thought has a perspective, not that time does \citep{Mellor1998}.
\end{reply}

\begin{objection}[Fragmentalism]
\citet{Fine2005} argues that McTaggart's argument can be resisted by denying that reality is coherent as a whole: tensed facts are each real, but not all together.
\end{objection}

\begin{reply}
This grants what the inversion needs. A reality whose tensed facts do not form a coherent whole is one in which the A-series is not coherent; and the inversion uses Assumption~\ref{a:m3} only to conclude that time does not consist in the A-series so understood.
\end{reply}

\section{Conclusion}

McTaggart's argument is valid, and the incoherence of the A-series is its strongest premise. Its weakest is the bridge, which owes its appeal to an equivocation on change. Run backwards from the reality of time, the argument refutes the bridge and leaves time to the B-series, as TSM-01 concludes; this is also the B-series as TSM-14 uses it, the causal and metrical structure of spacetime.

\begin{thebibliography}{9}
\bibitem[Dummett(1960)]{Dummett1960} Dummett, M. (1960). A defense of McTaggart's proof of the unreality of time. \emph{The Philosophical Review}, 69(4), 497--504.
\bibitem[Fine(2005)]{Fine2005} Fine, K. (2005). Tense and reality. In \emph{Modality and Tense: Philosophical Papers}, 261--320. Oxford University Press.
\bibitem[McTaggart(1908)]{McTaggart1908} McTaggart, J.~M.~E. (1908). The unreality of time. \emph{Mind}, 17(68), 457--474.
\bibitem[Mellor(1998)]{Mellor1998} Mellor, D.~H. (1998). \emph{Real Time II}. Routledge.
\bibitem[Russell(1903)]{Russell1903} Russell, B. (1903). \emph{The Principles of Mathematics}. Cambridge University Press.
\end{thebibliography}

\end{document}
